\section{Macro-level Production Function}
\label{sec:aggregation}

So far, the model has focused on individual firm decision-making: how a firm allocates production steps among automation, augmentation, and manual execution, and how its micro-level choices of task assignment and job design determine internal productivity and costs.
Ultimately, however, our interest extends beyond the firm to the broader economy. 
Improvements in AI quality not only alter how individual firms organize production but also affect the composition of inputs used across the economy and the aggregate relationship between labor, capital, and output.
Yet the firm-level optimization framework developed above prevents direct analysis of these aggregate effects and also makes our approach distinct from the large literature in which production is modeled directly at the task level and then aggregated to the economy level (e.g., \cite{acemoglu2022artificial, acemoglu2024tasks, acemoglu2025simple}).
To bridge this gap, this section examines whether the micro-level cost-minimization problem of firms can be aggregated into well-behaved firm- and economy-level production functions. 
In particular, we show that the detailed, many-input Leontief representation of production at the task level can be reduced to a more compact formulation with only a small number of effective inputs, and that aggregating across heterogeneous firms (in terms of AI quality level) yields a macro-level production function with a CES form.
This provides a tractable link between micro-level organization decisions and macroeconomic implications of improving AI quality.

We begin by showing how a firm's micro-level decisions can be explicitly represented by a Leontief production function at the task level.
Next, we demonstrate how, for the purposes of analyzing labor allocation within a firm, this detailed task-level production function can be simplified into a Leontief production function with just two aggregate inputs: skill-adjusted AI management labor and skill-adjusted manual labor (we will describe what we mean by ``skill-adjusted'' below).
We then explore how the production functions of numerous individual firms, which differ only in how they can effectively put the same AI technology into use, can be aggregated into a single, economy-wide CES production function.
In this macro-level representation, the inputs are economy-wide aggregates of the corresponding micro-level inputs from individual firms, thereby facilitating clearer insights into the overall labor allocation in the economy.
Finally, we put this representation to work and derive the general equilibrium effects of an improvement in AI quality on aggregate output, on the allocation of workers between AI management and manual work, and on wages.

Before the aggregation analysis, let us first describe the setup of production.
A firm uses two production inputs: (1) skill-adjusted AI management labor (to complete AI-assisted tasks), and (2) skill-adjusted manual labor (to perform manual tasks).
These two types of labor demand different compensations: executing an AI-assisted task with skill and time cost of one demands $w_{\AIletter}$, while completing a manual task with similar skill and time requirements commands a compensation of $w_{\manualLetter}$.
We call $w_{\AIletter}$ and $w_{\manualLetter}$ the \emph{base wage rates} for AI management and manual labor, respectively.

Firms use these two inputs to complete tasks.
Recall that the skill and time requirements of tasks are determined by the firm's AI strategy $\T$, while the aggregate skill requirements of jobs depend jointly on the AI strategy and job design $\J$.
Specifically, the per unit of time compensation of a worker employed to do tasks of a job is determined by all tasks in that job.
Consider Job 1 in Figure~\ref{fig:hierarchical_production} for example.
The worker assigned to perform Job 1 is required to obtain not only the skill required for manual Task 1, $\manualSkill{1}$, but must also possess the required skill for AI-assisted Task 2, $\AIskill{2}$, as the worker's total compensation per unit of time is determined at the job level, and equals $\manualSkill{1} + \AIskill{2}$.
Therefore, the more tasks firm includes in a job, the higher required compensation for every task in that job.

In order to produce, the firm hires manual labor to perform manual tasks and AI management labor for AI-assisted tasks.
Therefore, the total compensation of a job will be a weighted average of skill costs of the job's tasks, with weights being the base wage rates of each type of labor.
Specifically, the compensation for job $J$ will be:
\begin{equation}
w_{\manualLetter} \left(\sum_{T_b^{\manualLetter} \in J} \manualSkill{b}\right) + w_{\AIletter} \left(\sum_{T_b^{\AIletter} \in J} \AIskill{b}\right),\label{eq:new_wage}
\end{equation}
which is composed of the sum of contributions of skill costs from its manual tasks denoted by set $T_b^{\manualLetter}$ (each weighted by $w_{\manualLetter}$) and its AI-assisted tasks denoted by set $T_b^{\AIletter}$ (each weighted by $w_{\AIletter}$).\footnote{
This formulation is essentially equivalent to the original definition of wage in \eqref{eq:wage}, as pointed out in Footnote~\ref{foot:wage}. 
If we multiply each skill cost in \eqref{eq:new_wage} by its respective base wage rate and redefine the task's skill cost as the resulting product, we arrive at the original formulation given by \eqref{eq:wage}.
The formulation in \eqref{eq:new_wage} allows tasks to contribute differently to the job's wage depending on their mode of execution and type of labor that need to perform them.}

Next, we define skill-adjusted labor, which is the smallest unit of work in our analysis.
Let $J(b)$ denote the job to which task $b$ belongs, and $E(b)$ be the mode of execution of task $b$ (i.e., $E(b) = \manualLetter$ if $b$ is a manual task, or $E(b) = \AIletter$ if $b$ is an AI chain).
Define the skill-adjusted time requirement of task $b$ as
\[
\skillAdjustedTimeLetter_b = \frac{w_{\manualLetter} \left(\sum_{T_{\ell}^{\manualLetter} \in J(b)} \manualSkill{\ell}\right) + w_{\AIletter} \left(\sum_{T_{\ell}^{\AIletter} \in J(b)} \AIskill{\ell}\right)}{w_{E(b)}} \, \timecost{b}.\label{eq:skill_adjusted_time}
\]
Note that the fraction appearing behind $\timecost{b}$ represents an effective skill adjustment factor, as the numerator captures the total compensation of the job task $b$ belongs to, while the denominator normalizes by the base wage rate for the specific mode of execution of task $b$.
The resulting fraction thus has units of skill intensity.
This characterization becomes useful later as it allows expressing the (expected) wage bill paid for task $b$ simply as $w_{\AIletter}\,\skillAdjustedTimeLetter_{b}\,\alpha^{-d_b}$ if it is AI-assisted or as $w_{\manualLetter}\,\skillAdjustedTimeLetter_{b}$ if it is executed manually.


\subsection{Within-Firm Aggregation: Leontief to Leontief}
\label{sec:agg_within}

We now consider the firm's production.
To produce output, each firm must complete requirements of all \( m \) production steps \( s_1, \ldots, s_m \).
This naturally results in a Leontief production structure with multiple inputs.\footnote{
Different AI deployment strategies and job designs lead to different input proportions but share the same fundamental Leontief form.
}
Although firms choose whether to execute each step manually or by the help of AI, we show that analyzing labor allocation can be simplified.
Specifically, production function of the firm can be represented by a Leontief function with only two firm-level aggregate inputs: AI management labor and manual labor.

Suppose the firm solves the long-run cost minimization problem \eqref{eq:totalcost_with_handoff} for a given AI quality level $\alpha$, and obtains the optimal AI strategy \( \T \) and job design \( \J \).
Let \( |\T| = n \), implying the original \( m \) production steps are organized into \( n \) distinct tasks, and \( |\J| = p \), indicating these \( n \) tasks are grouped into \( p \) jobs.
Given the fixed AI strategy, the execution mode for each production step is predetermined. 
This lets us consider and work directly with tasks rather than individual steps.
Furthermore, with the job design held fixed, each firm takes the job boundaries and the total skill requirements within jobs as given.

We assume hand-offs are performed manually by the worker who completes the final task of each job.
This allows us to simplify the analysis in two ways:
(1) Since job design is fixed, we know exactly which steps occur at the job boundaries and thus incur hand-off time costs. Therefore, hand-off of job $J_j$ can be treated as a standalone, human-executed task with time requirement \( \handofftime{}(J_j) \) and skill requirement equal to the total skill requirement of job $J_j$.\footnote{
Denote the skill-adjusted time cost of hand-off task of job $J_j$ with $\skillAdjustedTimeHandoff{}(J_j)$.
}
(2) The order of tasks, including the newly defined hand-off tasks, can be rearranged within the production sequence.
Specifically, we relabel tasks such that tasks \( 1 \) to \( k \) are AI-assisted and thus require AI management labor, while tasks \( k+1 \) to \( n \), along with the hand-off tasks \( n+1 \) to \( n+p-1 \), are executed manually and require manual labor.\footnote{
Recall that the final job \( p \) requires no hand-off by definition. The last hand-off thus occurs at the end of job \( p-1 \).
}

The firm's task-level production can be expressed in the following Leontief form:
\begin{equation}
x = \min\left\{\frac{\labor{1}}{\skillAdjustedTimeAI{1}\,\alpha^{-d_{1}}},\,\cdots,\,\frac{\labor{k}}{\skillAdjustedTimeAI{k}\,\alpha^{-d_{k}}},\,\frac{\labor{k+1}}{\skillAdjustedTimeManual{k+1}},\,\cdots,\,\frac{\labor{n}}{\skillAdjustedTimeManual{n}},\,\frac{\labor{n+1}}{\skillAdjustedTimeHandoff{}(J_1)},\,\cdots,\,\frac{\labor{n+p-1}}{\skillAdjustedTimeHandoff{}(J_{p-1})}\right\}.\label{eq:task_level_prod}
\end{equation}
The \( x \) on the left-hand side represents the firm's output while $\labor{b}$ is the amount of labor assigned to task $b$. All other variables remain as previously defined.

Since the AI strategy and job design are fixed, the required amount of skill-adjusted time to spend on each task in the denominators of (\ref{eq:task_level_prod}) are fixed, and the firm takes them as given.
In equilibrium, allocation of labor to tasks satisfies
\[
x=\frac{\labor{1}}{\skillAdjustedTimeAI{1}\,\alpha^{-d_{1}}}=\cdots=\frac{\labor{k}}{\skillAdjustedTimeAI{k}\,\alpha^{-d_{k}}}=\frac{\labor{k+1}}{\skillAdjustedTimeManual{k+1}}=\cdots=\frac{\labor{n}}{\skillAdjustedTimeManual{n}}=\frac{\labor{n+1}}{\skillAdjustedTimeHandoff{}(J_1)}=\cdots=\frac{\labor{n+p-1}}{\skillAdjustedTimeHandoff{}(J_{p-1})}.
\]
Notice that the ratio of labor allocated to any two tasks $a$ and $b$ is independent of output level and wage rates (regardless of whether tasks $a$ and $b$ are executed manually or by the help of AI), and solely depends on their relative skill-adjusted time requirements, which are fixed given $\T$ and $\J$.
The fixed ratio of labor inputs implies that the rate of substitution between any pair of tasks is independent of the labor allocated to other tasks:
\begin{equation}
\frac{\partial}{\partial \labor{z}}\left(\frac{\frac{\Delta x}{\Delta \labor{a}}}{\frac{\Delta x}{\Delta \labor{b}}}\right)=0,\qquad\forall z\neq a,b.\label{eq:task_subst_rate_indep}
\end{equation}

A necessary and sufficient condition to aggregate the firm's production function from a Leontief with one input per task into a Leontief with (firm-level) aggregate AI management labor and manual labor is that the rate of substitution between any two tasks within the same aggregate input type is independent of all tasks in the other aggregate input type \citep{leontief1947introduction, fisher1965embodied, felipe2003aggregation}.
In other words, the relative amount of labor allocated to any two human-executed (AI-managed) tasks should depend exclusively on tasks within the human-executed (AI-managed) group.

The condition expressed in (\ref{eq:task_subst_rate_indep}) satisfies the aggregation requirement above.\footnote{In fact, this condition is stricter than necessary. It not only guarantees labor independence across aggregated input types, but also imposes a stronger condition: that the rate of substitution between tasks within the same aggregate input type is independent even of other tasks within the same aggregate type.}
Thus, the firm's production function can be represented as:
\begin{equation}
x=\min\,\left\{\frac{\bar{\alpha}\,\labor{\AIletter}}{\skillAdjustedTimeLetter_{\AIletter}},\,\frac{\labor{\manualLetter}}{\skillAdjustedTimeLetter_{\manualLetter}}\right\}.\label{eq:micro_agg_prod}
\end{equation}
Here, $\labor{\AIletter}$ and $\labor{\manualLetter}$ represent respectively the firm-level aggregate AI management labor and manual labor.
Moreover, the aggregate skill-adjusted AI and manual time requirements, $\skillAdjustedTimeLetter_{\AIletter}$ and $\skillAdjustedTimeLetter_{\manualLetter}$, are defined as the sum of skill-adjusted time requirements of tasks in their respective groups:
\begin{align}
\skillAdjustedTimeLetter_{\AIletter} & = \sum_{b=1}^{k}\skillAdjustedTimeAI{b}, \label{eq:def_t_AI}\\[8pt]
\skillAdjustedTimeLetter_{\manualLetter} & = \sum_{j=1}^{p-1}\skillAdjustedTimeHandoff{}(J_j) + \sum_{b=k+1}^{n}\skillAdjustedTimeManual{b}. \label{eq:def_t_H}
\end{align}
Finally, $\bar{\alpha}$ is the firm's effective AI quality level defined as
\begin{equation}
\bar{\alpha}=\frac{\sum_{b=1}^{k}\skillAdjustedTimeAI{b}}{\sum_{b=1}^{k}\skillAdjustedTimeAI{b}\alpha^{-d_{b}}},\label{eq:alpha_bar}
\end{equation}
so that the following relationship holds:
\[
\frac{\skillAdjustedTimeLetter_{\AIletter}}{\bar{\alpha}}=\sum_{b=1}^{k}\skillAdjustedTimeAI{b}\alpha^{-d_{b}}.
\]

In short, (\ref{eq:micro_agg_prod}) allows us to analyze labor allocation using the simpler two-input Leontief production function for a given AI strategy \( \T \) and job design \( \J \).
This firm-level aggregate Leontief function implies the equilibrium condition:
\begin{equation}
x=\frac{\bar{\alpha}\,\labor{\AIletter}}{\skillAdjustedTimeLetter_{\AIletter}}
=\,\frac{\labor{\manualLetter}}{\skillAdjustedTimeLetter_{\manualLetter}}.
\label{eq:micro_fixed_proportions}
\end{equation}
The representation in \eqref{eq:micro_agg_prod} helps us gain clearer insights into the allocation between manual and AI management labor at the firm level by removing the need to deal with the complex, many-input production function (\ref{eq:task_level_prod}) at the level of individual tasks.



\subsection{Cross-Firm Aggregation: Leontief to CES}

Next, we analyze whether Leontief production functions of many firms can be collectively represented by a single economy-wide production function, where each input in the macro function aggregates the corresponding inputs from the micro functions. 
To do so, we draw on the extensive literature on aggregation theorems in production theory, which establishes conditions for the existence of an aggregate production function based on individual firms' production functions.
The central idea is to introduce an appropriate form of heterogeneity across firms, allowing micro-level production functions to be smoothed into a well-behaved aggregate function (see \cite{sato1975production}, Chapters 2 and 4, for existence conditions.).\footnote{
For instance, \cite{houthakker1955pareto} studies aggregation from micro Leontief functions to a macro Cobb-Douglas function; \cite{levhari1968note} investigates aggregation from micro Leontief functions to a macro CES function; and \cite{sato1969micro} explores aggregation from micro CES to macro CES functions. 
See \cite{baqaee2019foundations} for a broader formulation of this problem.}
While most existence results in this literature do not yield a closed-form functional representation, we show here that under certain conditions on firm heterogeneity, firm-level Leontief production functions can be aggregated into an economy-wide CES production function.

Consider a unit mass of firms in the economy, each producing output with AI management labor, human labor, and capital.
We assume capital is a fixed input, so firms make labor decisions conditional on their given capital stock.
Moreover, we impose a Leontief structure between capital and labor inputs, normalizing the workflow so that exactly one unit of capital is required per unit of output produced.\footnote{This normalization can be interpreted as assuming identical capital productivity across firms. Since we introduce firm-level heterogeneity through variations in AI quality, we abstract away from additional sources of heterogeneity across firms.} 
Firms do not know their individual effective AI quality level before production occurs; instead, they share a common prior belief regarding the distribution from which these quality levels are drawn.

Production unfolds in two stages.
In the first stage, firms solve the long-run cost minimization problem \eqref{eq:totalcost_with_handoff}, committing to an AI deployment strategy $\T$ and a job design $\J$ based on their \emph{expected} AI quality levels.
Since firms hold the same expectations about the distribution of effective AI quality, they choose identical AI strategies and job designs.
In the second stage, firms learn their realized effective AI quality levels, hire the required labor and capital, and begin production.\footnote{As production depends on the realized effective AI quality level, $\bar{\alpha}$ also determines firms' sizes.}
The resulting Leontief production function, determined by the previously selected AI strategy, job design, and the firm's realized effective AI quality, takes a form similar to (\ref{eq:micro_agg_prod}).

This two-stage structure mirrors realistic firm behavior.
Firms initially decide how to structure operations and post job vacancies.
Only after making these structural commitments do they hire workers and acquire machines to begin production.
Once hiring occurs, firms discover the actual productivity levels of inputs, which necessitates operating under their previously chosen organizational structures.

In what follows, we demonstrate that the heterogeneity in effective AI quality level can be structured so that micro-level Leontief production functions (\ref{eq:micro_agg_prod}) aggregate into a macro-level CES production function, in which the CES inputs are aggregates of the micro-level inputs \citep{levhari1968note}.
Suppose specifically that the macro-level production function takes the form:
\begin{equation}
X = \left(\theta_{\AIletter}\,\aggregateAIlabor^{\rho}+\theta_{\manualLetter}\,\aggregateManualLabor^{\rho}+(1-\theta_{\AIletter}-\theta_{\manualLetter})\,K^\rho\right)^{\frac{1}{\rho}},\label{eq:macro_agg_prod}
\end{equation}
where $X$ represents the (economy-wide) aggregate output, $\aggregateAIlabor$ denotes aggregate AI management labor, $\aggregateManualLabor$ denotes aggregate manual labor, $K$ is aggregate capital, and parameters $\theta_{\AIletter}$ and $\theta_{\manualLetter}$ represent the weights of corresponding inputs.
Since the measure of firms in the economy is normalized to $1$, aggregate capital is also normalized to $1$, making the third term in (\ref{eq:macro_agg_prod}) effectively constant at $1-\theta_{\AIletter}-\theta_{\manualLetter}$.

To link the macro production function above to the micro production functions, we must relate aggregated firm-level inputs \( \labor{\AIletter}\) and \( \labor{\manualLetter} \) from (\ref{eq:micro_agg_prod}) to their corresponding economy-wide aggregates \( \aggregateAIlabor\) and \( \aggregateManualLabor \) in (\ref{eq:macro_agg_prod}).
This requires either determining the macro production function parameters from the distribution of heterogeneity, or specifying the form of heterogeneity given a set of macro production function parameters.
We adopt the latter approach and derive the output probability density function in terms of effective AI quality, denoted by \( \phi(\bar{\alpha}) \), as a function of \( \theta_{\AIletter}, \theta_{\manualLetter}, \rho \).
Throughout, we assume that $\rho<0$ (which implies elasticity of substitution $\sigma<1$), indicating that macro-level production exhibits some degree of complementarity between aggregate inputs.

Normalize the output price to $p=1$, so that $w_{\AIletter}$ and $w_{\manualLetter}$ can be interpreted as real wage rates for a unit of skill-adjusted AI management and manual labor, respectively.
A firm produces only if it earns nonnegative profits:
\[
w_{\AIletter}\,\labor{\AIletter}+w_{\manualLetter}\,\labor{\manualLetter}\leq x.
\]
Substituting for $\labor{\manualLetter}$ and $x$ from (\ref{eq:micro_fixed_proportions}) into the profitability condition yields:
\[
w_{\AIletter}\,\labor{\AIletter}+w_{\manualLetter}\,\frac{\skillAdjustedTimeLetter_{\manualLetter}\,\bar{\alpha}\,\labor{\AIletter}}{\skillAdjustedTimeLetter_{\AIletter}} \leq \frac{\bar{\alpha}\,\labor{\AIletter}}{\skillAdjustedTimeLetter_{\AIletter}}.
\]
Rearranging terms and canceling $\labor{\AIletter}$ gives the lower bound of firms' effective AI quality level:\footnote{We assume the parameters $w_{\AIletter},\,w_{\manualLetter},\,\skillAdjustedTimeLetter_{\AIletter},\,\skillAdjustedTimeLetter_{\manualLetter}$ are such that $0 < ({w_{\AIletter}\,\skillAdjustedTimeLetter_{\AIletter}})/(1-w_{\manualLetter}\,\skillAdjustedTimeLetter_{\manualLetter}) < 1$.}
\[
\bar{\alpha} \geq \frac{w_{\AIletter}\,\skillAdjustedTimeLetter_{\AIletter}}{1 - w_{\manualLetter}\,\skillAdjustedTimeLetter_{\manualLetter}}.
\]
This lower bound implies that only firms with realized effective AI quality level above this threshold will produce in equilibrium and others exit the market.
Moreover, observe from \eqref{eq:alpha_bar} that for $d_b \geq 0$ the upper bound of $\bar{\alpha}$ is 1 and is achieved when $\alpha \rightarrow 1^{-}$.
Therefore:
\begin{equation}
\frac{w_{\AIletter}\,\skillAdjustedTimeLetter_{\AIletter}}{1 - w_{\manualLetter}\,\skillAdjustedTimeLetter_{\manualLetter}} \leq \bar{\alpha} \leq 1.
\label{eq:alpha_threshold}
\end{equation}

Let $\phi(\bar{\alpha})$ be the distribution of output across firms according to their effective AI quality level. 
A firm with effective AI quality $\bar{\alpha}$ thus produces output $x = \phi(\bar{\alpha})$ by definition.
From equation (\ref{eq:micro_fixed_proportions}), it directly follows that the AI management labor and manual labor used by this firm are: $\labor{\AIletter} = (\skillAdjustedTimeLetter_{\AIletter}/\bar{\alpha})\,\phi(\bar{\alpha}),$ and $\labor{\manualLetter} = \skillAdjustedTimeLetter_{\manualLetter}\,\phi(\bar{\alpha}).$
Consequently, aggregate output $X$, aggregate AI management labor $\aggregateAIlabor$, and aggregate manual labor $\aggregateManualLabor$ can be expressed as:
\begin{align}
X & = \int_{\frac{w_{\AIletter}\,\skillAdjustedTimeLetter_{\AIletter}}{1 - w_{\manualLetter}\,\skillAdjustedTimeLetter_{\manualLetter}}}^{1}\,\phi(\bar{\alpha})\,d\bar{\alpha}, \label{eq:micro_X}\\
\aggregateAIlabor & = \int_{\frac{w_{\AIletter}\,\skillAdjustedTimeLetter_{\AIletter}}{1 - w_{\manualLetter}\,\skillAdjustedTimeLetter_{\manualLetter}}}^{1}\,\skillAdjustedTimeLetter_{\AIletter}\frac{\phi(\bar{\alpha})}{\bar{\alpha}}\,d\bar{\alpha}, \label{eq:micro_M}\\
\aggregateManualLabor & = \int_{\frac{w_{\AIletter}\,\skillAdjustedTimeLetter_{\AIletter}}{1 - w_{\manualLetter}\,\skillAdjustedTimeLetter_{\manualLetter}}}^{1}\,\skillAdjustedTimeLetter_{\manualLetter}\,\phi(\bar{\alpha})\,d\bar{\alpha}.\label{eq:micro_H}
\end{align}

Substituting the aggregated variables from the micro-level equations (\ref{eq:micro_X}, \ref{eq:micro_M}, \ref{eq:micro_H}) into the aggregate production function (\ref{eq:macro_agg_prod}), we obtain:
\begin{equation}
\int_{\frac{w_{\AIletter}\,\skillAdjustedTimeLetter_{\AIletter}}{1 - w_{\manualLetter}\,\skillAdjustedTimeLetter_{\manualLetter}}}^{1}\,\phi(\bar{\alpha})\,d\bar{\alpha}
=
\left(
\theta_{\AIletter}\left(\skillAdjustedTimeLetter_{\AIletter}\int_{\frac{w_{\AIletter}\,\skillAdjustedTimeLetter_{\AIletter}}{1 - w_{\manualLetter}\,\skillAdjustedTimeLetter_{\manualLetter}}}^{1}\,\frac{\phi(\bar{\alpha})}{\bar{\alpha}}\,d\bar{\alpha}\right)^{\rho}
+
\theta_{\manualLetter} \left(\skillAdjustedTimeLetter_{\manualLetter}\int_{\frac{w_{\AIletter}\,\skillAdjustedTimeLetter_{\AIletter}}{1 - w_{\manualLetter}\,\skillAdjustedTimeLetter_{\manualLetter}}}^{1}\,\phi(\bar{\alpha})\,d\bar{\alpha}\right)^{\rho}
+
\left(1-\theta_{\AIletter}-\theta_{\manualLetter}\right)
\right)^{\frac{1}{\rho}}.
\label{eq:macro_agg_prod_with_micro_aggregates}
\end{equation}
Solving for $\phi(\bar{\alpha})$ in terms of \( \theta_{\AIletter}, \theta_{\manualLetter}, \rho \) we get the following  distribution for effective AI quality:
\begin{equation}
\phi(\bar{\alpha})=\frac{\left(1-\theta_{\AIletter}-\theta_{\manualLetter}\right)^{\frac{1}{\rho}}\left(1-\theta_{\manualLetter}\,\skillAdjustedTimeLetter_{\manualLetter}^{\rho}\right)^{\frac{\rho}{\rho-1}}\left(\theta_{\AIletter}\,\skillAdjustedTimeLetter_{\AIletter}^{\rho}\right)^{\frac{1}{1-\rho}}}{1-\rho}(\bar{\alpha})^{\frac{1}{\rho-1}}\left[1-\theta_{\manualLetter}\,\skillAdjustedTimeLetter_{\manualLetter}^{\rho}-\left(1-\theta_{\manualLetter}\,\skillAdjustedTimeLetter_{\manualLetter}^{\rho}\right)^{\frac{\rho}{\rho-1}}\left(\theta_{\AIletter}\,\skillAdjustedTimeLetter_{\AIletter}^{\rho}\right)^{\frac{1}{1-\rho}}(\bar{\alpha})^{\frac{\rho}{\rho-1}}\right]^{-\frac{1+\rho}{\rho}}.\label{eq:phi}
\end{equation}
See Subsection~\ref{sec:aggregationDist} below for derivation details.

This completes the production function aggregation procedure.
The derived distribution (\ref{eq:phi}) characterizes firm-level heterogeneity in the effective AI quality level required to support a macro CES production function of the form (\ref{eq:macro_agg_prod}), with economy-wide aggregate AI management and manual labor inputs and parameters $\theta_{\AIletter}, \theta_{\manualLetter}, \rho$, obtained from micro-level Leontief production functions of the form (\ref{eq:micro_agg_prod}).


The aggregation results in this section provide a foundation for linking micro-level technology deployment and organizational choices to macroeconomic production frameworks.
By showing how heterogeneous firms with step-by-step Leontief technologies and endogenous task boundaries can be represented by a CES aggregator, our framework offers a micro-founded rationale for using aggregate CES production functions to study labor demand, substitution patterns, and productivity in economies adopting new AI technologies.
This connection clarifies how firm-level decisions about automation, augmentation, and job design scale up to shape aggregate technological change.



\subsection{General Equilibrium Effects of Improving AI Quality}
\label{sec:agg_compstat}

The point of an aggregate production function is to let us ask aggregate questions, and the most natural one here is what happens to output, employment, and wages as AI quality improves.
The difficulty is that better AI does not simply make a fixed set of AI-executed steps more reliable.
Its central effect in our model is to change the organization of production itself: as $\alpha$ rises, chains extend over steps that were previously performed by hand, and the boundary between what people do and what AI does moves.
Holding the AI strategy fixed would assume away the mechanism the model is about.
We therefore let the firm reoptimize, and show that the general equilibrium consequences of an improvement depend on the firm's entire response, reorganization included, only through how it moves two numbers.

\subsubsection*{Two Sufficient Statistics for the Firm's Response}

Let $(\T, \J)$ be the organization the firm chooses at AI quality $\alpha$, and define
\begin{equation}
T_{\AIletter}(\alpha) \;=\; \sum_{b=1}^{k} \skillAdjustedTimeAI{b}\,\alpha^{-d_{b}},
\qquad
T_{\manualLetter}(\alpha) \;=\; \skillAdjustedTimeLetter_{\manualLetter},
\label{eq:cs_requirements}
\end{equation}
where the first sum runs over the AI chains of $\T$, with $d_b$ the total difficulty of chain $b$, and $\skillAdjustedTimeLetter_{\manualLetter}$ is as in \eqref{eq:def_t_H}.
These are the two input coefficients of the firm's technology: $T_{\AIletter}$ is the skill-adjusted AI management labor and $T_{\manualLetter}$ the skill-adjusted manual labor that a best-practice firm needs per unit of output.
Once effective AI quality is decomposed as in Assumption~\ref{as:neutral_heterogeneity} below, the firm-level Leontief function \eqref{eq:micro_agg_prod} reads
\begin{equation}
x = \min\left\{\frac{\varepsilon\,\labor{\AIletter}}{T_{\AIletter}},\ \frac{\labor{\manualLetter}}{T_{\manualLetter}}\right\},
\label{eq:cs_micro_prod}
\end{equation}
and nothing else about the workflow, the chain structure, or the job design enters the aggregate economy.

An improvement in AI quality moves the pair $(T_{\AIletter}, T_{\manualLetter})$ through two distinct channels.
The first operates at a fixed organization: a higher $\alpha$ raises the success probability of every chain, so that fewer retries are needed and $T_{\AIletter}$ falls at rate
\begin{equation}
-\left.\frac{\partial \ln T_{\AIletter}}{\partial \ln \alpha}\right|_{\T}
\;=\;
\frac{\sum_{b=1}^{k}\skillAdjustedTimeAI{b}\,d_{b}\,\alpha^{-d_{b}}}{\sum_{b=1}^{k}\skillAdjustedTimeAI{b}\,\alpha^{-d_{b}}}
\;\equiv\; \bar{D} \;>\; 0,
\label{eq:cs_passthrough}
\end{equation}
the average difficulty of the firm's AI-executed work weighted by the expected time each chain consumes, while $T_{\manualLetter}$ does not move at all.
Because a chain's difficulty is the sum of the difficulties of its steps, $\bar{D}$ is larger when AI is deployed in longer chains, so chaining amplifies how much a given improvement in raw AI quality is worth.

The second is the chaining channel, and it is the one that makes the organization of production endogenous to $\alpha$.
When an improvement makes it worthwhile to extend a chain over a step that was previously manual, $T_{\manualLetter}$ falls by that step's manual time requirement, while $T_{\AIletter}$ \emph{rises}, because the lengthened chain carries a higher total difficulty and so must be retried more often.
Automation therefore appears in \eqref{eq:cs_requirements} as a substitution out of manual labor and into AI management labor at the level of the firm's input coefficients, and it is the only way $T_{\manualLetter}$ can move at all.

We write $\hat{T}_{\AIletter}$ and $\hat{T}_{\manualLetter}$ for the proportional changes in the two coefficients induced by a small improvement in AI quality, combining both channels, so that $\hat{T}_{\manualLetter} < 0$ exactly when the improvement pulls steps out of manual execution.
Because the firm chooses among finitely many organizations, the chaining channel operates at the discrete reorganization thresholds of Section~\ref{sec:nonlinear} rather than continuously, and $\hat{T}_{\manualLetter} = 0$ between thresholds.
The results below are stated for small changes in $(T_{\AIletter}, T_{\manualLetter})$.
They describe a threshold crossing to first order, and describe it exactly in the limit of a workflow with many steps, where thresholds are dense and each reorganization moves a small share of the work.

\subsubsection*{Deployment Heterogeneity}

The one substantive restriction we need concerns how a firm's effective AI quality relates to the technology it has deployed.

\begin{assumption}[Neutral deployment heterogeneity]
\label{as:neutral_heterogeneity}
A firm's effective AI quality factors as $\bar{\alpha} = \Lambda\,\varepsilon$, where $\Lambda$ is common across firms and depends only on the AI technology and the organization adopted, and $\varepsilon \in (0,1]$ is the firm's relative deployment effectiveness, whose distribution across productive capacity is invariant to both.
\end{assumption}

This holds exactly if firms differ in an \emph{integration reliability} $\varepsilon$ that multiplies the success probability of every AI chain the firm runs, so that a chain spanning steps $(s_\ell, \dotsc, s_r)$ succeeds with probability $\varepsilon \prod_{i=\ell}^{r} q_i$ rather than $\prod_{i=\ell}^{r} q_i$.
This captures how well a firm has wired AI into its workflow, as distinct from how good the technology itself is.
Substituting into \eqref{eq:alpha_bar} gives $\bar{\alpha} = \varepsilon\,\Lambda$ with $\Lambda = \skillAdjustedTimeLetter_{\AIletter}/T_{\AIletter}$, the effective AI quality of a firm that deploys the technology at best practice, so that the firm's AI management requirement per unit of output is $\skillAdjustedTimeLetter_{\AIletter}/\bar{\alpha} = T_{\AIletter}/\varepsilon$, as in \eqref{eq:cs_micro_prod}.
This is the precise sense in which the firms considered above ``differ only in how they can effectively put the same AI technology into use'': $\Lambda$ is the technology as deployed and $\varepsilon$ is the firm's ability to deploy it.
Requiring the distribution of $\varepsilon$ to be invariant is natural here because $\phi$ describes how productive capacity, equivalently the fixed capital stock, is distributed across firms of differing deployment effectiveness, and neither the frontier technology nor the firm's choice of workflow redistributes capital across firms.

With this decomposition the aggregation of the previous subsection goes through with $\varepsilon$ in place of $\bar{\alpha}$ as the index of firm heterogeneity, and it does so for \emph{any} organization the firm might adopt.
What is more, the organization affects the resulting CES only through its weights.

\begin{lem}[Invariance of the aggregate elasticity of substitution]
\label{lem:sigma_invariance}
Under Assumption~\ref{as:neutral_heterogeneity}, for any input coefficients $(T_{\AIletter}, T_{\manualLetter})$ the economy with heterogeneity distribution \eqref{eq:phi} admits the CES representation
\begin{equation}
X = \left(\theta'_{\AIletter}\,\aggregateAIlabor^{\rho}+\theta'_{\manualLetter}\,\aggregateManualLabor^{\rho}+\theta'_{K}\,K^\rho\right)^{\frac{1}{\rho}},
\label{eq:macro_agg_prod_general}
\end{equation}
with the same $\rho$, and hence the same elasticity of substitution $\sigma = 1/(1-\rho)$, and with weights $(\theta'_{\AIletter}, \theta'_{\manualLetter}, \theta'_{K})$ that depend on the organization and reduce to $(\theta_{\AIletter}, \theta_{\manualLetter}, 1-\theta_{\AIletter}-\theta_{\manualLetter})$ at the baseline organization.
A reorganization of the workflow therefore shifts the weights of the macro production function but leaves its elasticity of substitution unchanged.
\end{lem}

\begin{proof}
Under Assumption~\ref{as:neutral_heterogeneity} the firm's technology is \eqref{eq:cs_micro_prod} and the break-even condition \eqref{eq:alpha_threshold} becomes $\varepsilon \geq u \equiv w_{\AIletter}T_{\AIletter}/(1 - w_{\manualLetter}T_{\manualLetter})$, which are the expressions of Section~\ref{sec:aggregation} with $\varepsilon$ in place of $\bar{\alpha}$ and $(T_{\AIletter}, T_{\manualLetter})$ in place of $(\skillAdjustedTimeLetter_{\AIletter}, \skillAdjustedTimeLetter_{\manualLetter})$.
Hence $X = \Gamma(u)$, $\aggregateAIlabor = T_{\AIletter}\Psi(u)$ and $\aggregateManualLabor = T_{\manualLetter}\Gamma(u)$, with $\Gamma$ and $\Psi$ as in \eqref{eq:gamma} and \eqref{eq:psi}.
From Section~\ref{sec:aggregationDist}, $\Gamma(u) = \theta_{K}^{1/\rho}\left(b_0 - b_1 v\right)^{-1/\rho}$ and $\Psi(u) = (b_1/b_0)\,u^{1/(\rho-1)}\,\Gamma(u)$, where $v \equiv u^{\rho/(\rho-1)}$, $\theta_K = 1-\theta_{\AIletter}-\theta_{\manualLetter}$, $b_0 = 1 - \theta_{\manualLetter}\skillAdjustedTimeLetter_{\manualLetter}^{\rho}$ and $b_1 = b_0^{\rho/(\rho-1)}(\theta_{\AIletter}\skillAdjustedTimeLetter_{\AIletter}^{\rho})^{1/(1-\rho)}$.
Substituting into \eqref{eq:macro_agg_prod_general} and dividing by $X^{\rho}$ gives
\[
1 \;=\; \theta'_{\AIletter}\,T_{\AIletter}^{\rho}\left(\frac{b_1}{b_0}\right)^{\rho} v \;+\; \theta'_{\manualLetter}\,T_{\manualLetter}^{\rho} \;+\; \frac{\theta'_{K}}{\theta_{K}}\left(b_0 - b_1 v\right),
\]
which holds identically in $u$ if and only if the coefficient on $v$ vanishes and the constant term equals one.
Together with $\theta'_{\AIletter}+\theta'_{\manualLetter}+\theta'_{K} = 1$ this is a system of three linear equations in the three weights with a unique solution, which is the baseline solution when $(T_{\AIletter}, T_{\manualLetter}) = (\skillAdjustedTimeLetter_{\AIletter}, \skillAdjustedTimeLetter_{\manualLetter})$.\footnote{
The weights remain positive for organizations in a neighborhood of the baseline by continuity.
Consistency of the closed form \eqref{eq:phi} with the definition $\Gamma(u) = \int_u^1\phi(\varepsilon)\,d\varepsilon$ also requires the boundary condition $\Gamma(1) = 0$, which holds if and only if $\theta_{\AIletter}\skillAdjustedTimeLetter_{\AIletter}^{\rho} + \theta_{\manualLetter}\skillAdjustedTimeLetter_{\manualLetter}^{\rho} = 1$.
We impose this restriction on the baseline macro parameters throughout, which costs one degree of freedom in $(\theta_{\AIletter}, \theta_{\manualLetter}, \rho)$ and leaves the aggregation result otherwise unchanged.
}
\end{proof}

\subsubsection*{Closing the Model}

To move from technology to equilibrium we must say how the two labor inputs are supplied.
Workers are ex ante identical and are hired and trained by firms for designated jobs (Section~\ref{sec:jobs_wages}), so free entry into both kinds of work equalizes the base wage rates, $w_{\AIletter} = w_{\manualLetter} = w$, which is the normalization the main text adopts in Footnote~\ref{foot:wage}.
Total skill-adjusted labor is supplied inelastically at $L$, and capital remains the fixed factor, in supply $K = 1$.

A firm with deployment effectiveness $\varepsilon$ produces $\phi(\varepsilon)$ if it breaks even and nothing otherwise, so it produces exactly when $\varepsilon \geq u$, where the entry cutoff $u$ and the wage satisfy
\begin{equation}
u \;=\; \frac{w\,T_{\AIletter}}{1 - w\,T_{\manualLetter}},
\qquad\text{equivalently}\qquad
w \;=\; \left(\frac{T_{\AIletter}}{u} + T_{\manualLetter}\right)^{-1}.
\label{eq:cs_cutoff}
\end{equation}
The second expression says the wage is the inverse of the labor the marginal firm needs per unit of output.
Aggregate output and the two labor aggregates are
\begin{equation}
X = \Gamma(u), \qquad \aggregateAIlabor = T_{\AIletter}\,\Psi(u), \qquad \aggregateManualLabor = T_{\manualLetter}\,\Gamma(u),
\label{eq:cs_aggregates}
\end{equation}
and the labor market clears when
\begin{equation}
T_{\AIletter}\,\Psi(u) + T_{\manualLetter}\,\Gamma(u)
\;=\;
\int_u^1 \left(\frac{T_{\AIletter}}{\varepsilon} + T_{\manualLetter}\right)\phi(\varepsilon)\,d\varepsilon
\;=\; L.
\label{eq:cs_clearing}
\end{equation}
The middle expression is the economy's total labor requirement, the labor each producing firm needs per unit of output integrated against its output.
It is strictly decreasing in $u$, vanishes as $u \to 1$, and diverges as $u \to 0$, so \eqref{eq:cs_clearing} has a unique solution $u \in (0,1)$ for any $L > 0$, and \eqref{eq:cs_cutoff} then delivers the wage.\footnote{
Divergence at $u \to 0$ follows because $\phi(\varepsilon)$ behaves like $\varepsilon^{1/(\rho-1)}$ near zero by \eqref{eq:phi}, so $\phi(\varepsilon)/\varepsilon$ is not integrable at the origin.
Note also that firms with $\varepsilon < u$ leave their capital idle, so aggregate capital in use, $\Gamma(u)$, is below the aggregate supply $K = 1$.
}
Firms above the cutoff earn rents on their capacity, and we write
\[
s_{\AIletter} = \frac{w\,\aggregateAIlabor}{X},
\qquad
s_{\manualLetter} = \frac{w\,\aggregateManualLabor}{X} = w\,T_{\manualLetter},
\qquad
s_{K} = 1 - s_{\AIletter} - s_{\manualLetter}
\]
for the shares of output accruing to AI management labor, to manual labor, and to capital.

\subsubsection*{Comparative Statics}

Proposition~\ref{prop:agg_compstat} gives the general equilibrium response to an improvement in AI quality, allowing the firm to reorganize.

\begin{proposition}[General equilibrium effects of AI progress]
\label{prop:agg_compstat}
Suppose an improvement in AI quality changes the firm's input coefficients by $(\hat{T}_{\AIletter}, \hat{T}_{\manualLetter})$, and let
\begin{equation}
\gamma \;\equiv\; -\left(s_{\AIletter}\hat{T}_{\AIletter} + s_{\manualLetter}\hat{T}_{\manualLetter}\right)
\label{eq:cs_tfp}
\end{equation}
denote the induced rate of aggregate productivity growth.
Then, under Assumption~\ref{as:neutral_heterogeneity},
\begin{align}
\hat{X} \;=\; \gamma,
\qquad
\hat{u} \;=\; -\frac{s_{K}}{\sigma\,s_{\AIletter}}\,\gamma,
\qquad
&\hat{\aggregateManualLabor} \;=\; \gamma + \hat{T}_{\manualLetter},
\qquad
\hat{\aggregateAIlabor} \;=\; -\frac{\aggregateManualLabor}{\aggregateAIlabor}\left(\gamma + \hat{T}_{\manualLetter}\right),
\label{eq:cs_quantities}\\[4pt]
\hat{w} \;=\; \underbrace{\frac{1-s_{\manualLetter}}{s_{\AIletter}}\left(1 - \frac{s_{K}}{\sigma}\right)\gamma}_{\text{productivity}}
\;&+\;
\underbrace{\frac{s_{K}\,s_{\manualLetter}}{s_{\AIletter}}\,\hat{T}_{\manualLetter}}_{\text{displacement}}.
\label{eq:cs_prices}
\end{align}
\end{proposition}

\begin{proof}
Let $\eta \equiv u\,\phi(u)/X = -\,d\ln\Gamma/d\ln u$ denote the \emph{entry elasticity}, the elasticity of aggregate output with respect to the entry cutoff.
We proceed in four steps.

\emph{Step 1: the entry margin.}
Write the left-hand side of \eqref{eq:cs_clearing} as $g(u, T_{\AIletter}, T_{\manualLetter})$.
Since $\Gamma'(u) = -\phi(u)$ and $\Psi'(u) = -\phi(u)/u$,
\[
\frac{\partial g}{\partial \ln T_{\AIletter}} = \aggregateAIlabor,
\qquad
\frac{\partial g}{\partial \ln T_{\manualLetter}} = \aggregateManualLabor,
\qquad
\frac{\partial g}{\partial u} = -\phi(u)\left(\frac{T_{\AIletter}}{u} + T_{\manualLetter}\right) = -\frac{\phi(u)}{w},
\]
where the last equality uses \eqref{eq:cs_cutoff}.
Totally differentiating \eqref{eq:cs_clearing} therefore gives
\begin{equation}
\hat{u} \;=\; \frac{w\left(\aggregateAIlabor \hat{T}_{\AIletter} + \aggregateManualLabor \hat{T}_{\manualLetter}\right)}{u\,\phi(u)} \;=\; \frac{s_{\AIletter}\hat{T}_{\AIletter} + s_{\manualLetter}\hat{T}_{\manualLetter}}{\eta} \;=\; -\frac{\gamma}{\eta}.
\label{eq:cs_proof_cutoff}
\end{equation}

\emph{Step 2: quantities.}
From \eqref{eq:cs_aggregates}, $\hat{X} = -\eta\,\hat{u} = \gamma$, and $\aggregateManualLabor = T_{\manualLetter}X$ gives $\hat{\aggregateManualLabor} = \hat{T}_{\manualLetter} + \gamma$.
Market clearing $\aggregateAIlabor = L - \aggregateManualLabor$ then implies $\aggregateAIlabor\,\hat{\aggregateAIlabor} = -\aggregateManualLabor\,\hat{\aggregateManualLabor}$, which is the last expression in \eqref{eq:cs_quantities}.

\emph{Step 3: the wage.}
Differentiating the second expression in \eqref{eq:cs_cutoff} and using $w\,T_{\AIletter}/u = 1 - w\,T_{\manualLetter} = 1 - s_{\manualLetter}$,
\[
\hat{w} \;=\; -(1-s_{\manualLetter})\left(\hat{T}_{\AIletter} - \hat{u}\right) - s_{\manualLetter}\hat{T}_{\manualLetter}.
\]
Substituting $\hat{u} = -\gamma/\eta$ and $\hat{T}_{\AIletter} = -(\gamma + s_{\manualLetter}\hat{T}_{\manualLetter})/s_{\AIletter}$ from \eqref{eq:cs_tfp}, and collecting terms,
\[
\hat{w} \;=\; \frac{1-s_{\manualLetter}}{s_{\AIletter}}\left(1 - \frac{s_{\AIletter}}{\eta}\right)\gamma \;+\; \left(\frac{1-s_{\manualLetter}}{s_{\AIletter}} - 1\right)s_{\manualLetter}\hat{T}_{\manualLetter},
\]
and $(1-s_{\manualLetter})/s_{\AIletter} - 1 = s_{K}/s_{\AIletter}$.

\emph{Step 4: the entry elasticity.}
It remains to show that $\eta = \sigma\,s_{\AIletter}/s_{K}$ for any $(T_{\AIletter}, T_{\manualLetter})$.
Using $\Psi(u) = (b_1/b_0)\,u^{1/(\rho-1)}\Gamma(u)$ from the proof of Lemma~\ref{lem:sigma_invariance} together with $w\,T_{\AIletter} = u\,(1 - w\,T_{\manualLetter}) = u\,(1-s_{\manualLetter})$,
\[
s_{\AIletter} = \frac{w\,T_{\AIletter}\Psi(u)}{\Gamma(u)} = w\,T_{\AIletter}\,\frac{b_1}{b_0}\,u^{\frac{1}{\rho-1}} = (1-s_{\manualLetter})\,\frac{b_1 v}{b_0},
\qquad v \equiv u^{\frac{\rho}{\rho-1}},
\]
so that $s_{K} = 1 - s_{\manualLetter} - s_{\AIletter} = (1-s_{\manualLetter})(b_0 - b_1 v)/b_0$ and therefore $s_{\AIletter}/s_{K} = b_1 v/(b_0 - b_1 v)$.
Since $\Gamma(u) = \theta_{K}^{1/\rho}(b_0 - b_1 v)^{-1/\rho}$ and $d\ln v/d\ln u = \rho/(\rho-1)$,
\[
\eta = -\frac{d \ln \Gamma}{d \ln u} = \frac{1}{\rho}\,\frac{d \ln (b_0 - b_1 v)}{d \ln u} = \frac{b_1 v}{(1-\rho)(b_0 - b_1 v)} = \sigma\,\frac{s_{\AIletter}}{s_{K}}.
\]
Note that $\eta$ depends on the organization only through the equilibrium cutoff it induces.
Substituting into \eqref{eq:cs_proof_cutoff} and into Step~3 gives \eqref{eq:cs_quantities} and \eqref{eq:cs_prices}.
\end{proof}

The two channels through which AI progress reaches the economy have sharply different aggregate signatures, and Corollaries~\ref{cor:cs_within} and~\ref{cor:cs_reorg} record them.

\begin{cor}[Better AI at a fixed organization]
\label{cor:cs_within}
Between reorganization thresholds, $\hat{T}_{\manualLetter} = 0$ and $\hat{T}_{\AIletter} = -\bar{D}\,\hat{\alpha}$, so $\gamma = \bar{D}\,s_{\AIletter}\,\hat{\alpha}$ and
\[
\hat{X} = \hat{\aggregateManualLabor} = \bar{D}\,s_{\AIletter}\,\hat{\alpha} > 0,
\qquad
\hat{\aggregateAIlabor} = -\,\bar{D}\,s_{\manualLetter}\,\hat{\alpha} < 0,
\qquad
\hat{w} = \bar{D}\left(1 - s_{\manualLetter}\right)\left(1 - \frac{s_{K}}{\sigma}\right)\hat{\alpha}.
\]
Output and manual employment rise, AI management employment falls, and the wage rises if and only if $\sigma > s_{K}$.
\end{cor}

\begin{cor}[Reorganization that lengthens a chain]
\label{cor:cs_reorg}
At a reorganization that pulls a manual step into a chain, $\hat{T}_{\AIletter} > 0 > \hat{T}_{\manualLetter}$, and $\gamma > 0$ whenever the reorganization lowers the economy's labor requirement.
Then $\hat{X} > 0$, $\hat{u} < 0$, $\hat{\aggregateManualLabor} < 0$ and $\hat{\aggregateAIlabor} > 0$.
Output and entry rise, manual employment falls, AI management employment rises, and the wage carries a strictly negative displacement term in \eqref{eq:cs_prices}.
\end{cor}

\begin{proof}
Only the employment signs require an argument.
By \eqref{eq:cs_tfp}, $\hat{\aggregateManualLabor} = \gamma + \hat{T}_{\manualLetter} = -s_{\AIletter}\hat{T}_{\AIletter} + (1-s_{\manualLetter})\hat{T}_{\manualLetter}$, and both terms are negative when $\hat{T}_{\AIletter} > 0 > \hat{T}_{\manualLetter}$.
The sign of $\hat{\aggregateAIlabor}$ follows from market clearing.
\end{proof}

Four features of these results are worth drawing out.

First, aggregate output grows at exactly the rate at which the firm's technology improves.
The object $\gamma$ in \eqref{eq:cs_tfp} is the standard productivity residual, the share-weighted reduction in unit input requirements, and with labor and capital in fixed supply it is also the growth rate of aggregate output and of output per worker.
Within a regime it equals $\bar{D}\,s_{\AIletter}$, so the macroeconomic return to better AI is governed by how much of the wage bill is currently spent managing AI and by how difficult the work AI has been given is, rather than by how impressive the improvement is in engineering terms.
Since $\gamma = \bar{D}\,s_{\AIletter}$ between thresholds, and $\bar{D}$ is non-increasing in $\alpha$ within a regime but jumps up whenever a reorganization lengthens a chain, the non-monotonicity of Lemma~\ref{lem:reorg_threshold} carries over to aggregate productivity growth: it diminishes as a given organization of work is pushed further and revives when the workflow is redrawn, which is the sense in which our framework microfounds a productivity J-curve \citep{brynjolfsson2021productivity}.\footnote{
That $\bar{D}$ falls within a regime follows because it is a weighted average of chain difficulties with weights proportional to $\skillAdjustedTimeAI{b}\alpha^{-d_b}$, and a higher $\alpha$ shifts those weights toward the easier chains; formally $d\bar{D}/d\ln\alpha$ equals minus the variance of $d_b$ under those weights.
When a reorganization absorbs a manual step into chain $b$, both $d_b$ and its weight rise, so $\bar{D}$ jumps up.
}

Second, and this is where the chaining margin does its work, the two channels move employment composition in \emph{opposite} directions.
At a fixed organization, each unit of output still requires $T_{\manualLetter}$ units of manual labor while the human time needed per unit of AI-executed work falls, so the economy reallocates workers out of AI management and into manual work, and the manual steps that AI has not absorbed become the binding constraint on scale.
A reorganization reverses this: it removes manual steps from the workflow outright, and by Corollary~\ref{cor:cs_reorg} manual employment falls and AI management employment rises even though output expands.
Whether AI progress shrinks or grows manual employment over a stretch of the diffusion path is therefore a question about how often, and how far, firms cross reorganization thresholds along the way, not a question about the rate of technical improvement as such.
Without the chaining margin the first channel is all there is, and better AI could only ever push workers toward manual work.

Third, the wage response in \eqref{eq:cs_prices} is a race between the labor AI releases and the output expansion that reabsorbs it, in the sense of \cite{acemoglu2019automation}, and the decomposition separates the two forces cleanly.
The productivity term is positive if and only if $\sigma > s_{K}$.
Behind that condition is the entry elasticity $\eta = u\phi(u)/X = \sigma\,s_{\AIletter}/s_{K}$, which measures how much productive capacity sits just below the break-even level of deployment effectiveness: when $\eta$ is large, a small decline in the break-even level brings a great deal of output on line and labor demand expands.
The displacement term is active only when the improvement reorganizes work, and it is then unambiguously negative.
Automation therefore puts downward pressure on wages over and above anything that happens to productivity, and it can push the wage down even in an economy where $\sigma > s_{K}$ and pure quality gains would raise it.

Fourth, the elasticity of substitution that governs all of this is a property of the dispersion of deployment effectiveness across firms, not of the workflow.
By Lemma~\ref{lem:sigma_invariance}, when firms reorganize around better AI the macro production function's weights move, shifting toward AI management labor, but $\sigma$ does not.
Aggregate substitutability between AI management labor and the rest of the economy is thus inherited from how unevenly firms deploy AI, and improvements in AI quality change what the economy produces and who produces it without changing how substitutable the inputs are.



\subsection{Derivation of the Effective AI Quality Heterogeneity Distribution}
\label{sec:aggregationDist}

In this subsection we describe how to obtain the analytical formula for distribution of heterogeneity $\phi(\bar{\alpha})$ expressed in (\ref{eq:phi}) by extending the method of \cite{levhari1968note}.

Define $u=({w_{\AIletter}\,\skillAdjustedTimeLetter_{\AIletter}})/(1-w_{\manualLetter}\,\skillAdjustedTimeLetter_{\manualLetter})$ and rewrite equation (\ref{eq:macro_agg_prod_with_micro_aggregates}) as:
\begin{equation}
\left(\int_{u}^{1}\,\phi(\bar{\alpha})\,d\bar{\alpha}\right)^{\rho}
=
\theta_{\AIletter}\left(\skillAdjustedTimeLetter_{\AIletter}\int_{u}^{1}\,\frac{\phi(\bar{\alpha})}{\bar{\alpha}}\,d\bar{\alpha}\right)^{\rho}
+
\theta_{\manualLetter} \left(\skillAdjustedTimeLetter_{\manualLetter}\int_{u}^{1}\,\phi(\bar{\alpha})\,d\bar{\alpha}\right)^{\rho}
+
\left(1-\theta_{\AIletter}-\theta_{\manualLetter}\right).
\label{eq:macro_agg_prod_with_hetDist}
\end{equation}
Our goal is to solve for $\phi(\bar{\alpha})$ in the equation above as a function of $\theta_{\AIletter},\theta_{\manualLetter},\rho$.
Define:
\begin{equation}
\Gamma(u)=\int_{u}^{1}\,\phi(\bar{\alpha})\,d\bar{\alpha},\label{eq:gamma}
\end{equation}
and
\begin{equation}
\Psi(u)=\int_{u}^{1}\,\frac{\phi(\bar{\alpha})}{\bar{\alpha}}\,d\bar{\alpha}.\label{eq:psi}
\end{equation}
Note that $\Gamma^{'}(u)=-\phi(u)$ and $\Psi^{'}(u)=-\phi(u)/u$.
Rewrite equation (\ref{eq:macro_agg_prod_with_hetDist}) as:
\[
\Gamma^{\rho}(u)=\theta_{\AIletter}\,\skillAdjustedTimeLetter_{\AIletter}^{\rho}\,\Psi^{\rho}(u)+\theta_{\manualLetter}\,\skillAdjustedTimeLetter_{\manualLetter}^{\rho}\,\Gamma^{\rho}(u)+(1-\theta_{\AIletter}-\theta_{\manualLetter}).
\]
Rearranging terms, we obtain:
\begin{equation}
\left(1-\theta_{\manualLetter}\,\skillAdjustedTimeLetter_{\manualLetter}^{\rho}\right)\Gamma^{\rho}(u)=\theta_{\AIletter}\,\skillAdjustedTimeLetter_{\AIletter}^{\rho}\,\Psi^{\rho}(u)+(1-\theta_{\AIletter}-\theta_{\manualLetter}).\label{eq:macro_agg_pro_with_gamma}
\end{equation}

Differentiating equation (\ref{eq:macro_agg_pro_with_gamma}) with respect to $u$ yields:
\[
\left(1-\theta_{\manualLetter}\,\skillAdjustedTimeLetter_{\manualLetter}^{\rho}\right)\Gamma^{\rho-1}(u)=\theta_{\AIletter}\,\skillAdjustedTimeLetter_{\AIletter}^{\rho}\,\Psi^{\rho-1}(u)\,u^{-1}.
\]
Solving for $\Psi^{\rho-1}(u)$ as a function of $u$ and $\Gamma^{\rho-1}(u)$ we get:
\[
\Psi^{\rho-1}(u)=\frac{\left(1-\theta_{\manualLetter}\,\skillAdjustedTimeLetter_{\manualLetter}^{\rho}\right)}{\theta_{\AIletter}\,\skillAdjustedTimeLetter_{\AIletter}^{\rho}}\Gamma^{\rho-1}(u)\,u.
\]
Finally, express $\Psi(u)$ in terms of $\Gamma(u)$:
\begin{equation}
\Psi(u)=\left(1-\theta_{\manualLetter}\,\skillAdjustedTimeLetter_{\manualLetter}^{\rho}\right)^{\frac{1}{\rho-1}}\left(\theta_{\AIletter}\,\skillAdjustedTimeLetter_{\AIletter}^{\rho}\right)^{\frac{1}{1-\rho}}u^{\frac{1}{\rho-1}}\,\Gamma(u).\label{eq:psi_solved}
\end{equation}


Substituting \( \Psi(u) \) from equation (\ref{eq:psi_solved}) into equation (\ref{eq:macro_agg_pro_with_gamma}), we have:
\[
\left(1-\theta_{\manualLetter}\,\skillAdjustedTimeLetter_{\manualLetter}^{\rho}\right)\Gamma^{\rho}(u)=\left(1-\theta_{\manualLetter}\,\skillAdjustedTimeLetter_{\manualLetter}^{\rho}\right)^{\frac{\rho}{\rho-1}}\left(\theta_{\AIletter}\,\skillAdjustedTimeLetter_{\AIletter}^{\rho}\right)^{\frac{1}{1-\rho}}u^{\frac{\rho}{\rho-1}}\,\Gamma^{\rho}(u)+\left(1-\theta_{\AIletter}-\theta_{\manualLetter}\right).
\]
Rearranging terms to solve for \( \Gamma(u) \), we get:
\[
\Gamma(u)=\left(1-\theta_{\AIletter}-\theta_{\manualLetter}\right)^{\frac{1}{\rho}}\left[1-\theta_{\manualLetter}\,\skillAdjustedTimeLetter_{\manualLetter}^{\rho}-\left(1-\theta_{\manualLetter}\,\skillAdjustedTimeLetter_{\manualLetter}^{\rho}\right)^{\frac{\rho}{\rho-1}}\left(\theta_{\AIletter}\,\skillAdjustedTimeLetter_{\AIletter}^{\rho}\right)^{\frac{1}{1-\rho}}u^{\frac{\rho}{\rho-1}}\right]^{-\frac{1}{\rho}}.
\]

Recall that \( \Gamma'(u)=-\phi(u) \).
Differentiating \( \Gamma(u) \), we obtain:
\begin{align*}
\Gamma'(u)&=\frac{\partial}{\partial u}\left[\left(1-\theta_{\AIletter}-\theta_{\manualLetter}\right)^{\frac{1}{\rho}}\left[1-\theta_{\manualLetter}\,\skillAdjustedTimeLetter_{\manualLetter}^{\rho}-\left(1-\theta_{\manualLetter}\,\skillAdjustedTimeLetter_{\manualLetter}^{\rho}\right)^{\frac{\rho}{\rho-1}}\left(\theta_{\AIletter}\,\skillAdjustedTimeLetter_{\AIletter}^{\rho}\right)^{\frac{1}{1-\rho}}u^{\frac{\rho}{\rho-1}}\right]^{-\frac{1}{\rho}}\right] \\
&=\frac{\left(1-\theta_{\AIletter}-\theta_{\manualLetter}\right)^{\frac{1}{\rho}}\left(1-\theta_{\manualLetter}\,\skillAdjustedTimeLetter_{\manualLetter}^{\rho}\right)^{\frac{\rho}{\rho-1}}\left(\theta_{\AIletter}\,\skillAdjustedTimeLetter_{\AIletter}^{\rho}\right)^{\frac{1}{1-\rho}}}{\rho-1}u^{\frac{1}{\rho-1}}\left[1-\theta_{\manualLetter}\,\skillAdjustedTimeLetter_{\manualLetter}^{\rho}-\left(1-\theta_{\manualLetter}\,\skillAdjustedTimeLetter_{\manualLetter}^{\rho}\right)^{\frac{\rho}{\rho-1}}\left(\theta_{\AIletter}\,\skillAdjustedTimeLetter_{\AIletter}^{\rho}\right)^{\frac{1}{1-\rho}}u^{\frac{\rho}{\rho-1}}\right]^{-\frac{1+\rho}{\rho}}.
\end{align*}
Thus, we have:
\begin{equation}
\phi(\bar{\alpha})=\frac{\left(1-\theta_{\AIletter}-\theta_{\manualLetter}\right)^{\frac{1}{\rho}}\left(1-\theta_{\manualLetter}\,\skillAdjustedTimeLetter_{\manualLetter}^{\rho}\right)^{\frac{\rho}{\rho-1}}\left(\theta_{\AIletter}\,\skillAdjustedTimeLetter_{\AIletter}^{\rho}\right)^{\frac{1}{1-\rho}}}{1-\rho}(\bar{\alpha})^{\frac{1}{\rho-1}}\left[1-\theta_{\manualLetter}\,\skillAdjustedTimeLetter_{\manualLetter}^{\rho}-\left(1-\theta_{\manualLetter}\,\skillAdjustedTimeLetter_{\manualLetter}^{\rho}\right)^{\frac{\rho}{\rho-1}}\left(\theta_{\AIletter}\,\skillAdjustedTimeLetter_{\AIletter}^{\rho}\right)^{\frac{1}{1-\rho}}(\bar{\alpha})^{\frac{\rho}{\rho-1}}\right]^{-\frac{1+\rho}{\rho}},
\end{equation}
which is the formula provided in (\ref{eq:phi}) above.
